\documentclass[a4paper,12pt]{article}
\usepackage[utf8]{inputenc}
\usepackage[T1]{fontenc}
\usepackage[ngerman]{babel}
\usepackage{amsmath}
\usepackage{amssymb}
\usepackage{enumitem}
\usepackage{geometry}
\geometry{a4paper, margin=2.5cm}
\usepackage{parskip}

\title{Einsendeaufgaben -- Aufgabe 5.1\\[0.5em]
\large Modul 61112: Lineare Algebra}
\author{}
\date{}

\begin{document}
\maketitle

\section*{Aufgabe 1}

\begin{enumerate}[label=\alph*)]

    \item
    \begin{align*}
        \chi_A &= \det(XI_3 - A) \\
               &= \det\begin{pmatrix}
                    X-a & -b  & c \\
                    -b  & X-a & 0 \\
                    -c  & 0   & X-a
                  \end{pmatrix} \\
               &= (X-a)^3 - (X-a)b^2 - (X-a)c^2 \\
               &= (X-a)\bigl((X-a)^2 - b^2 - c^2\bigr) \\
               &= (X-a)(X^2 - 2aX + a^2 - b^2 - c^2)
    \end{align*}

    \textit{Fall 1: $a = b = c = 0$}

    $A$ ist eine Nullmatrix und demnach eine Diagonalmatrix.

    \textit{Fall 2: $a \neq 0$, $b = c = 0$}

    Es folgt $\chi_A = (X-a)^3$. Wegen
    \begin{align*}
        a(a) = 3 &= \dim(\operatorname{Ker}(O_3)) \\
                 &= \dim(\operatorname{Ker}(A - aI_3)) \\
                 &= \gamma(a)
    \end{align*}
    ist $A$ nach Satz 5.3.20 diagonalisierbar.

    \textit{Fall 3: $b \neq 0$ oder $c \neq 0$}

    Die Eigenwerte sind
    \[
        \lambda_1 = a \quad\text{und}\quad \lambda_{2/3} = a \pm \sqrt{b^2 + c^2}.
    \]
    Wegen $b \neq 0$ oder $c \neq 0$ ist $\lambda_1 \neq \lambda_2 \neq \lambda_3$.
    Demnach existieren drei Eigenwerte.

    Da $\sum_{i=1}^{3} \gamma(\lambda_i) \leq \dim(U)$, muss
    \[
        \gamma(\lambda) = 1 = a(\lambda) \quad\text{für jedes } \lambda \in \{\lambda_1, \lambda_2, \lambda_3\}
    \]
    gelten. Nach Satz 5.3.20 ist also $A$ diagonalisierbar.

    \item
    Nach a) Fall 3 folgt
    \[
        \lambda_1 = 1 \quad\text{und}\quad \lambda_{2/3} = 1 \pm 7.
    \]

    \begin{align*}
        \operatorname{Eig}_A(1) &= \operatorname{Ker}(A - I_3) \\
            &= \operatorname{Ker}\begin{pmatrix}
                0 & 0 & 7 \\
                0 & 0 & 0 \\
                7 & 0 & 0
              \end{pmatrix} \\
            &= \left\{ \lambda \begin{pmatrix} 0 \\ 1 \\ 0 \end{pmatrix} \;\middle|\; \lambda \in \mathbb{C} \right\}
    \end{align*}

    \begin{align*}
        \operatorname{Eig}_A(1+7) &= \operatorname{Ker}(A - (1+7)I_3) \\
            &= \operatorname{Ker}\begin{pmatrix}
                -7 & 0  & 7 \\
                0  & -7 & 0 \\
                7  & 0  & -7
              \end{pmatrix} \\
            &= \operatorname{Ker}\begin{pmatrix}
                1 & 0 & -1 \\
                0 & 1 & 0 \\
                0 & 0 & 0
              \end{pmatrix} \\
            &= \left\{ \lambda \begin{pmatrix} 1 \\ 0 \\ 1 \end{pmatrix} \;\middle|\; \lambda \in \mathbb{C} \right\}
    \end{align*}

    \begin{align*}
        \operatorname{Eig}_A(1-7) &= \operatorname{Ker}(A - (1-7)I_3) \\
            &= \operatorname{Ker}\begin{pmatrix}
                7 & 0 & 7 \\
                0 & 7 & 0 \\
                7 & 0 & 7
              \end{pmatrix} \\
            &= \operatorname{Ker}\begin{pmatrix}
                1 & 0 & 1 \\
                0 & 1 & 0 \\
                0 & 0 & 0
              \end{pmatrix} \\
            &= \left\{ \lambda \begin{pmatrix} 1 \\ 0 \\ -1 \end{pmatrix} \;\middle|\; \lambda \in \mathbb{C} \right\}
    \end{align*}

    \[
        S := \begin{pmatrix}
            1 & 0 & 1 \\
            0 & 1 & 0 \\
            1 & 0 & -1
          \end{pmatrix}
    \]

    \begin{align*}
        &\left(\begin{array}{ccc|ccc}
            1 & 0 & 1 & 1 & 0 & 0 \\
            0 & 1 & 0 & 0 & 1 & 0 \\
            1 & 0 & -1 & 0 & 0 & 1
          \end{array}\right) \\
        \rightsquigarrow\;&\left(\begin{array}{ccc|ccc}
            1 & 0 & 1 & 1 & 0 & 0 \\
            0 & 1 & 0 & 0 & 1 & 0 \\
            0 & 0 & -2 & -1 & 0 & 1
          \end{array}\right) \\
        \rightsquigarrow\;&\left(\begin{array}{ccc|ccc}
            1 & 0 & 1 & 1 & 0 & 0 \\
            0 & 1 & 0 & 0 & 1 & 0 \\
            0 & 0 & 1 & \tfrac{1}{2} & 0 & -\tfrac{1}{2}
          \end{array}\right) \\
        \rightsquigarrow\;&\left(\begin{array}{ccc|ccc}
            1 & 0 & 0 & \tfrac{1}{2} & 0 & \tfrac{1}{2} \\
            0 & 1 & 0 & 0 & 1 & 0 \\
            0 & 0 & 1 & \tfrac{1}{2} & 0 & -\tfrac{1}{2}
          \end{array}\right)
    \end{align*}

    \[
        \Rightarrow\; S^{-1} = \begin{pmatrix}
            \tfrac{1}{2} & 0 & \tfrac{1}{2} \\
            0 & 1 & 0 \\
            \tfrac{1}{2} & 0 & -\tfrac{1}{2}
          \end{pmatrix}
    \]

    \begin{align*}
        S^{-1} A S &= S^{-1} \begin{pmatrix}
                1 & 0 & 7 \\
                0 & 1 & 0 \\
                7 & 0 & 1
              \end{pmatrix}\begin{pmatrix}
                1 & 0 & 1 \\
                0 & 1 & 0 \\
                1 & 0 & -1
              \end{pmatrix} \\
            &= \begin{pmatrix}
                \tfrac{1}{2} & 0 & \tfrac{1}{2} \\
                0 & 1 & 0 \\
                \tfrac{1}{2} & 0 & -\tfrac{1}{2}
              \end{pmatrix}\begin{pmatrix}
                8 & 0 & -6 \\
                0 & 1 & 0 \\
                8 & 0 & 6
              \end{pmatrix} \\
            &= \begin{pmatrix}
                8 & 0 & 0 \\
                0 & 1 & 0 \\
                0 & 0 & -6
              \end{pmatrix}
            = \begin{pmatrix}
                1+7 & 0 & 0 \\
                0 & 1 & 0 \\
                0 & 0 & 1-7
              \end{pmatrix}
    \end{align*}

    \item
    \begin{align*}
        \chi_A &= \det(XI_3 - A) \\
               &= \det\begin{pmatrix}
                    X  & -1 & -i \\
                    -1 & X  & 0 \\
                    -i & 0  & X
                  \end{pmatrix} \\
               &= X(X^2 - 1 + 1) \\
               &= X^3
    \end{align*}

    \[
        \dim(\operatorname{Eig}_{\chi_A}(0)) = \dim\left(\operatorname{Ker}\begin{pmatrix}
            0 & 1 & i \\
            1 & 0 & 0 \\
            i & 0 & 0
          \end{pmatrix}\right) = 1
    \]

    Demnach $\gamma(0) < a(0)$. Nach Satz 5.3.20 ist $A$ nicht diagonalisierbar.

\end{enumerate}

\end{document}
